\subsection{* Convex envelope of a weakly compact set}

THEOREM 3 (Krein). --- Let E be a Hausdorff and quasi-complete locally convex space, and let T be a topology on E compatible with the duality between E and E'. Let A be a subset of E which is compact for T. Then the closed convex balanced envelope C of A is compact for T.

We shall first make several reductions.

\begin{enumerate}
\def\labelenumi{\Alph{enumi})}
\item
  The set C is precompact for \(\mathcal{T}\) (II, p.~25, prop. 3), and A is compact for \(\sigma(\mathrm{E}, \mathrm{E}')\) . On account of prop. 1 (IV, p.~32), it is enough to prove that C is compact for \(\sigma(\mathrm{E}, \mathrm{E}')\) , and so we have reduced to the case where \(\mathcal{T} = \sigma(\mathrm{E}, \mathrm{E}')\) .
\item
  Since C is precompact and closed for \(\sigma(\mathrm{E}, \mathrm{E}')\) , it is bounded and closed for the initial topology of E (III, p.~3, prop. 2 and IV, p.~1, prop. 1); hence it is complete since E is quasi-complete. In other words, C is the closed convex balanced envelope of A in the completion \(\hat{\mathbf{E}}\) of E. Since the topology \(\sigma (\hat{\mathbf{E}},\mathrm{E}^{\prime})\) induces \(\sigma (\mathrm{E},\mathrm{E}^{\prime})\) on E, we have reduced to the case when E is complete.
\item
  Let \(\Gamma\) be the convex balanced envelope of A. Then C is the closure of \(\Gamma\) for \(\sigma(E, E')\) . By Eberlein's theorem (IV, p.~35, th. 1), it is enough to prove that every sequence \((x_n)_{n \in \mathbf{N}}\) of points of \(\Gamma\) has a limit point for \(\sigma(E, E')\) in E. But \(x_n\) belongs to the convex balanced envelope of a finite subset \(\mathbf{B}_n\) of A. Let F be the closed vector subspace of E generated by the countable set \(\mathbf{B} = \bigcup_{n} \mathbf{B}_n\) . Then F is complete, the topology \(\sigma(F, F')\) on F is induced by \(\sigma(E, E')\) and we have \(x_n \in F\) for all \(n \in \mathbf{N}\) . Hence it is enough to prove that \((x_n)_{n \in \mathbf{N}}\) has a limit point for \(\sigma(F, F')\) , which gives the reduction to the case when there exists a countable dense set in E.
\end{enumerate}

Let A be assigned the topology induced by \(\sigma(E, E')\) , which makes it a compact space. We define a linear mapping \(u: E' \to \mathcal{C}(A)\) by

\[
u (x ^ {\prime}) (a) = \langle a, x ^ {\prime} \rangle \quad (a \in \mathrm{A}, x ^ {\prime} \in \mathrm{E} ^ {\prime}).\tag{7}
\]

Let \((x_{n}^{\prime})_{n\in\mathbb{N}}\) be an equicontinuous sequence in E', converging to 0 for \(\sigma(\mathrm{E}^{\prime},\mathrm{E})\) . Then the sequence of functions \(u(x_{n}^{\prime})\) is bounded in \(\mathcal{C}(\mathrm{A})\) and converges simply to 0. For every \(\mu\in\mathcal{C}(\mathrm{A})^{\prime}\) , we have \(\lim_{n\to\infty}\mu(u(x_{n}^{\prime}))=0\) by lemma 2 (IV, p.~37). By the criterion given in the remark in III, p.~21, the linear form \(\mu\circ u\) on E' is then continuous for \(\sigma(\mathrm{E}^{\prime},\mathrm{E})\) for every \(\mu\in\mathcal{C}(\mathrm{A})^{\prime}\) . Hence there exists a linear mapping \(v:\mathcal{C}(\mathrm{A})^{\prime}\to\mathrm{E}\) satisfying the relation

\[
\langle u (x ^ {\prime}), \mu \rangle = \langle v (\mu), x ^ {\prime} \rangle \quad \left(x ^ {\prime} \in \mathrm{E} ^ {\prime}, \mu \in \mathcal {C} (\mathrm{A}) ^ {\prime}\right).\tag{8}
\]

It is clear that \(v\) is continuous if \(\mathcal{C}(\mathbf{A})'\) is assigned the topology \(\sigma(\mathcal{C}(\mathbf{A})', \mathcal{C}(\mathbf{A}))\) and \(E\) the topology \(\sigma(E, E')\) .

The unit ball (closed) B of the Banach space \(\mathcal{C}(\mathrm{A})\) is compact for the topology \(\sigma(\mathcal{C}(\mathrm{A})', \mathcal{C}(\mathrm{A}))\) (III, p.~17, cor. 3). Consequently, \(v(\mathrm{B})\) is a convex balanced and compact subset of E for \(\sigma(\mathrm{E}, \mathrm{E}')\) . For every \(a \in A\) , the continuous linear form \(\varepsilon_a : f \mapsto f(a)\) on \(\mathcal{C}(\mathrm{A})\) belongs to B, and we have \(v(\varepsilon_a) = a\) by formulas (7) and (8). Hence, \(A \subset v(\mathrm{B})\) , and so \(C \subset v(\mathrm{B})\) . This proves that C is compact for \(\sigma(\mathrm{E}, \mathrm{E}')\) .

Q.E.D.

